\section{The question and the calculations}
A star of $0.1\Ms$ burns hydrogen for several trillion years. Laughlin,
Bodenheimer and Adams \cite{lba97} (LBA97) followed its evolution into
white-dwarf cooling with a grey atmosphere and the opacities and reaction
rates of the 1990s. The aim of Ember is to follow such
a star from formation through hydrogen exhaustion, the transition to a helium
white dwarf, cooling through 1000, 500 and 100 K, and the conditional stages
that follow: encounters, accretion, dark-matter heating, and disappearance if
nucleons decay. The remit is to use the simplest model that resolves each
question and to keep numerical uncertainty below about a tenth of the
physical variation that matters.

The paper reports five calculations. Their clocks differ and they have not
been joined into one track.
\begin{itemize}
\item The \textbf{continuous calculation} starts from a Hayashi-track model
and runs deuterium, proton--proton and CN burning, plasma neutrino losses and
convective mixing in one program. With hot microscopic transport and
composition-dependent atmospheres it reaches $\continuousAge\Tyr$.
Age is measured from the contracting starting model (Section~\ref{sec:contraction}).
\item \textbf{Track S} (settling) starts from a static hydrogen-burning model
at age zero and follows hydrogen, helium-3, carbon, nitrogen and a moving
metal group with microscopic diffusion and settling. It reaches
$4.002\Tyr$ and supplies the initial state of the flash comparisons.
\item \textbf{Track F} (fixed metals) uses the same network and diffusion
with the metal abundance held fixed; it reaches $3.704\Tyr$.
\item \textbf{Track P} (proton--proton only) omits carbon--nitrogen burning
and microscopic diffusion; it reaches $3.890\Tyr$ and is the basis of the
LBA97 comparison.
\item The \textbf{flash continuations} F1 and F2 start from Track S after
its outer boundary is replaced by a pure-hydrogen white-dwarf table. They are
provisional (Section~\ref{sec:flash}).
\end{itemize}

Sections~\ref{sec:model}--\ref{sec:hydrogen} give the model, the
contraction results and the hydrogen-burning results. Section~\ref{sec:flash}
separates what is established about the helium-3 flash from what is not.
Section~\ref{sec:comparisons} compares with LBA97, with MESA and with
observed stars. Section~\ref{sec:limits} lists the missing physics that can
change the interpretation, Section~\ref{sec:lifetime} gives the conditional
lifetime timeline, and Section~\ref{sec:computing} records cost and
reproducibility. The appendix holds reusable implementation detail.
Every number in the text is a calculated value unless labelled as a
literature value or a conditional estimate. Table~\ref{tab:status}
summarizes the status of each calculation.

\begin{table}[H]\centering\small
\caption{Status of the calculations reported here.}\label{tab:status}
\begin{tabular}{@{}>{\raggedright\arraybackslash}p{0.19\textwidth}>{\raggedright\arraybackslash}p{0.18\textwidth}>{\raggedright\arraybackslash}p{0.26\textwidth}>{\raggedright\arraybackslash}p{0.28\textwidth}@{}}\toprule
Calculation & Clock & Last computed state & Status \\\midrule
Continuous calculation & from Hayashi model & $\continuousAge\Tyr$, \continuousTemperature\ K, surface $X=\continuousHydrogen$ & radiative core, convective envelope \\
Fixed-metal atmosphere comparison & from Hayashi model & $3.078\Tyr$, 3130 K & complete comparison; fully convective \\
Track S & static H-burning model & $4.002\Tyr$, 4612 K & complete to atmosphere limit \\
Track F & static H-burning model & $3.704\Tyr$, 3702 K & complete (comparison) \\
Track P & static H-burning model & $3.890\Tyr$, 5546 K & complete (comparison) \\
Flash continuations F1, F2 & from post-switch model & 5.744 Myr; 42.48 Myr & provisional; no continuous flash \\
Continuous settling calculation & from Hayashi model & below 4600 K & further cooling needs opacity and atmosphere coverage \\\bottomrule
\end{tabular}
\end{table}
